%% 302_Elementarzellen_Spektrale_Antwort_De_ch.tex
%% FFGFT -- Dok. 302: Elementarzellen, Ein-Bit-Kapazität und Partitionsinvarianz;
%%                    eine spektrale Antwort (Informationseinheit = Windungsquant Delta w = 1)
%% Autor: Johann Pascher
%% Datum: Juli 2026
%% Sprache: Deutsch

\section*{Dok.~302 \quad Elementarzellen, Ein-Bit-Kapazität und Partitionsinvarianz ---
eine spektrale Antwort}

Der folgende Text beantwortet einen naheliegenden Einwand gegen jedes Bild eines
\enquote{quantisierten Raums}: Welches beobachterunabhängige Kriterium unterscheidet
eine Elementarzelle von ihren Nachbarn? Warum sollte eine fundamentale Länge feste
Zellen mit je genau einem Bit implizieren? Ist die daraus berechnete Entropie
invariant unter Änderungen von Ursprung, Orientierung, Koordinaten oder Foliation?
Und folgt die Unteilbarkeit einer Zelle aus der Planck-Skala, oder ist sie eine
zusätzliche Annahme? Die kurze Antwort lautet: Jeder dieser Punkte ist korrekt --- als
Einwand gegen eine \emph{Zell-Ontologie} ---, und FFGFT beruht auf keiner solchen. Es
ersetzt Zellen durch topologische und spektrale Invarianten (Windungsklassen und Moden), und die Einwände lösen sich der Reihe nach auf.

\section{Das Bild, auf das die Einwände zielen}

Eine Ontologie \enquote{ein Bit pro Planck-Zelle} kachelt den Raum (oder eine
Horizontfläche) in feste Zellen der Größe ${\sim}\,\ell_P$, weist jeder eine
Ein-Bit-Kapazität zu und liest die Entropie aus der Zellzahl ab. Gegen dieses Bild
sind die folgenden Punkte stichhaltig:

\begin{itemize}
\item es gibt kein beobachterunabhängiges Kriterium, das eine Raumzelle von ihrer
Nachbarin unterscheidet: eine Kachelung trägt einen gewählten Ursprung und eine
gewählte Orientierung;
\item die Planck-Länge ist eine \emph{Skala}, und eine Skala impliziert für sich
genommen keine Zerlegung in diskrete Behälter;
\item eine partitionsbasierte Entropie ist nicht invariant unter Änderung von
Ursprung, Orientierung, Koordinaten oder Foliation --- sie erbt die Willkür der
Partition;
\item ist der Raum ein Kontinuum, so lässt sich jede Zelle weiter teilen, sodass
\enquote{eine Zelle $=$ ein unteilbares Bit} einen unabhängigen Grund braucht, warum
Halbzellen-Zustände verboten sind;
\item und ohne diesen Grund ist die Zelle schlicht ein zusätzliches Postulat, das
auf der Planck-Skala aufsitzt.
\end{itemize}

Alle fünf werden zugestanden. Genau deshalb beruht das Modell nicht auf Raumzellen.

\section{Invarianten statt Zellen}

Das fundamentale Objekt ist kein Gitter aus Zellen, sondern eine kompakte, randlose
Mannigfaltigkeit --- ein Vier-Torus $T^4$ (mit einer $\mathbb{Z}_3$-Identifikation),
eingebettet in einen Hilbert-Raum $\mathcal{H}=L^2(T^4)\otimes\mathbb{C}^3$. Die
elementaren diskreten Objekte der Theorie sind keine Raumbereiche; sie sind
\emph{Windungsklassen} ($\pi_1(T^4)\cong\mathbb{Z}^4$, die Träger der
Informationseinheit) und \emph{Eigenmoden} eines Operators auf dieser Mannigfaltigkeit
(der Verbindungs-Laplace- bzw.\ Massenoperator, die Träger des Massenspektrums). Mit
diesem einen Wechsel der Ontologie --- Invarianten statt Zellen --- lösen sich die
Punkte der Reihe nach.

\begin{enumerate}
\item \textbf{Beobachterunabhängige Individuierung.} Die elementaren diskreten Objekte
werden durch \emph{Invarianten} unterschieden, nicht durch eine Lage im Raum: die
Informationseinheit durch ihre \emph{Windungszahl} (eine topologische Invariante,
$\pi_1(T^4)$), das Massenspektrum durch \emph{Eigenwerte} (eine spektrale Invariante).
Beide hängen nicht von Koordinaten, Ursprung oder Orientierung ab. Es gibt keine
\enquote{Nachbarzelle im Raum} zu unterscheiden; es gibt diskrete, durch Invarianten
gekennzeichnete Klassen. Die Individuierung ist topologisch bzw.\ spektral, nicht
räumlich --- und genau das macht sie beobachterunabhängig.

\item \textbf{Was die Informationseinheit ist --- und was am \enquote{einen Bit}
Definition bleibt.} Die Diskretheit stammt nicht von der Planck-Länge und ist keine
Kachelung, sondern ist \emph{topologisch}: $\pi_1(T^4) \cong \mathbb{Z}^4$ kennt nur
ganzzahlige Windungszahlen. Die minimale Informationseinheit ist daher \emph{per
Definition} das Windungsquant $\Delta w = 1$ (Dok.~257) --- skalenuniversell und
energieunabhängig; ein halbes Windungsquant gibt es nicht. Die Identifikation
\enquote{ein Bit $\;\widehat{=}\;$ eine minimale Windung} ist eine \emph{Definition},
wenn auch eine topologisch motivierte (die kleinste nichttriviale Windung), keine
willkürliche Zwei-Zustands-Annahme. Das $\ln 2$ ist \emph{nicht} die fundamentale
Einheit, sondern ihr thermodynamischer Schatten: die \emph{Energie} eines Bits ist
skalenspezifisch, $E_\text{bit}(L) = \hbar c / L$, und Landauers $E_L = k_B T \ln 2$
ist der Spezialfall bei einer einzigen (der biologischen) Skala. Damit ist Takayukis
Punkt~(2) beantwortet: \enquote{warum genau ein Bit} ist eine Definition auf
topologischer Grundlage ($\Delta w = 1$), nicht aus einer Skala erzwungen --- und die
Einheit ist skalenuniversell, also gerade nicht an eine gewählte Partition gebunden.

\item \textbf{Invarianz unter Ursprung, Orientierung, Koordinaten, Foliation.} Ja. Da
die Entropie eine Funktion des Spektrums ist --- eine Zählung oder Spur über Eigenwerte
--- und Eigenwerte Invarianten sind, ist die Entropie automatisch unabhängig von
Ursprung, Orientierung, Koordinatenwahl und Basis. Dies ist kein zusätzliches Axiom,
sondern eine Eigenschaft spektraler Größen. Eine partitionsbasierte Entropie fiele
durch den Invarianztest --- womit der Einwand recht behält ---, weshalb das Modell die
Entropie gar nicht über eine Partition definiert. Ein weiteres, modellspezifisches
Resultat macht die Ursprung-/Orientierungs-Unabhängigkeit explizit: auf der kompakten
Struktur gibt es keine ausgezeichnete Indexzuweisung --- alle $D_6$-äquivalenten
Segmente der Torusstruktur sind gleichwertige Beschreibungen desselben Objekts
(\enquote{Verzahnung, nicht Ausschnitt}). Eine Auszeichnung einer \enquote{ersten}
Zelle oder einer bevorzugten Achse ist weder verfügbar noch nötig.

\item \textbf{Kontinuierlicher Raum, keine Halbzellen.} Der Raum ist kontinuierlich:
$T^4$ ist eine glatte Mannigfaltigkeit, ohne Raumatomismus. Man kann Unteilbarkeit
daher nicht durch Unterteilung des Raums gewinnen, und das Modell versucht es nicht.
Die Unteilbarkeit ist \emph{topologisch und spektral}, nicht räumlich:
$\pi_1(T^4)\cong\mathbb{Z}^4$ lässt keine halbe Windung zu, und das Spektrum auf der
kompakten Mannigfaltigkeit hat einen kleinsten, durch die Topologie festgelegten
Abstand (Umfang / Randlosigkeit), so wie eine endliche Saite nicht bei \enquote{der
Hälfte ihrer Grundfrequenz} schwingen kann. \enquote{Halbzellen-Zustände} ist in
dieser Sprache ein Kategorienfehler: es gibt keine Zellen zu halbieren. Das Kontinuum
der Mannigfaltigkeit und die Diskretheit des Spektrums bestehen spannungsfrei
nebeneinander --- genau wie in der Akustik, wo ein kontinuierliches Medium eine diskrete
Menge von Normalmoden trägt.

\item \textbf{Abgeleitet, nicht angenommen.} Die Diskretheit folgt; sie ist kein
zusätzliches \enquote{Zellpostulat}. Auf einer kompakten Mannigfaltigkeit ist die
Diskretheit des Spektrums ein Satz. Angenommen wird die kompakte Geometrie selbst ---
der Torus $T^4/\mathbb{Z}_3$ und der Wert von $\xi$. Die ehrliche Bilanz lautet: das
\enquote{Ein-Bit-pro-Zelle}-Postulat wird gegen ein \enquote{Kompakte-Geometrie}-Postulat
eingetauscht, und aus letzterem ist die Diskretheit abgeleitet, nicht injiziert. Das
ist eine echte Annahme, und dort sollte das Modell angegriffen werden --- aber
empirisch, nicht per Dekret: dieselbe Geometrie, die das diskrete Spektrum liefert,
legt auch Teilchenmassen fest, und das sind falsifizierbare Zahlen (etwa eine
$\tau$-Masse von $1776{,}969~\mathrm{MeV}$, prüfbar an Belle~II). Ist die Geometrie
falsch, fallen Spektrum und Massen gemeinsam.
\end{enumerate}

\section{Was angenommen, abgeleitet und gesetzt ist}

Auch die Geometrie selbst ist eine \emph{Beschreibung} (Dok.~301, zweite Ebene) und als
solche notwendig unvollständig; keine Beschreibung rekonstruiert das Objekt vollständig.
Fundamental im vollen Sinn ist allein das vollständige Objekt (dritte Ebene). Dass die
resonanzgetreue Geometrie ein \emph{reales} Objekt nachspurt, ist daher eine
ausdrückliche ontologische \emph{Setzung} --- keine Ableitung und nicht selbst
falsifizierbar; falsifizierbar sind allein ihre resonanzgetreuen Folgen (die Massen).
Diese Setzung wird benannt, nicht als abgeleitet ausgegeben.

Fundamental ist die Informationseinheit als \emph{Windungsquant} $\Delta w = 1$
(skalenuniversell, Dok.~257); \emph{Energie} dagegen ist eine skalenspezifische
\emph{Projektion} der Geometrie, $E_\text{bit}(L) = \hbar c / L$. Eine \emph{Fläche}
teilt diesen projizierten, skalenabhängigen Charakter: der Vier-Torus ist randlos
($\partial T^4 = \emptyset$) und trägt keine Horizontfläche, an der Bits säßen; eine
Fläche ist dort eine \emph{Projektionsfläche}, ein Objekt der zweiten Ebene. Wo eine
Entropie-Flächen-Relation $S = N\ln 2$ mit $N \propto A/\xi^2$ auftritt, zählt $N$
Windungsquanten und $\ln 2$ ist der thermodynamische Schatten; das $A$ ist keine
fundamentale Raumpartition (womit Takayukis Zell-Sorge wiederkäme), sondern eine
gewählte Projektion. Fundamental (dritte Ebene) ist weder Fläche noch Energie, sondern
das vollständige Objekt; die skalenuniverselle Einheit auf ihm ist $\Delta w = 1$.

Die ehrliche Bilanz lautet somit: \emph{deklarierte Annahme} --- die kompakte Geometrie
($T^4/\mathbb{Z}_3$, $\xi$); \emph{Definition} --- die Identifikation \enquote{Bit
$\widehat{=}\,\Delta w = 1$}, topologisch begründet; \emph{abgeleitet} --- die
Diskretheit ($\pi_1(T^4)\cong\mathbb{Z}^4$), die skalenspezifische Bit-Energie
$E_\text{bit} = \hbar c/L$ und die Invarianz; \emph{empirische Exposition} --- die
Massen; \emph{benannte Setzung} --- die Realität des Objekts.

\section{Zusammenfassung}

Die Einwände sind gegen eine Raumzell-Ontologie korrekt, und das Modell stimmt ihnen
zu, indem es keine ist. Die Diskretheit ist topologisch ($\pi_1(T^4)\cong\mathbb{Z}^4$)
und spektral (das Spektrum auf einer kompakten Mannigfaltigkeit); die minimale
Informationseinheit ist das Windungsquant $\Delta w = 1$ (das \enquote{Bit}, per
Definition, topologisch begründet), nicht die Kapazität einer Raumzelle; und die
Invarianz gilt, weil Windungszahlen und Spektren Invarianten sind, nicht weil eine
Partition sorgfältig gewählt wurde.

\section{Wo dies ausgearbeitet ist}

Das diskrete Spektrum des $T^4$-Verbindungs-Laplace-Operators wird in Dok.~283
behandelt; die Hilbert-Raum-Formulierung $\mathcal{H}=L^2(T^4)\otimes\mathbb{C}^3$ mit
Massen als Eigenwerten in Dok.~282; die Informationseinheit als Windungsquant
$\Delta w = 1$ und die Drei-Ebenen-Hierarchie (topologisch / geometrisch /
thermodynamisch) in Dok.~257 (auf Dok.~255 aufbauend); die Ursprung-/Orientierungs-Un\-abhängigkeit
(\enquote{Verzahnung, nicht Ausschnitt}) in Dok.~300. Insbesondere der Invarianzpunkt
verdient eine ausführlichere schriftliche Behandlung, als er bislang erhalten hat.

